\documentclass[11pt]{article}
\usepackage[a4paper,margin=18mm]{geometry}
\usepackage[no-math]{fontspec}
\setmainfont{Latin Modern Roman}
\usepackage{amsmath,amssymb,amsthm,xfrac,xspace,hyperref,graphicx,comment}
\excludecomment{DefTralics}

\providecommand{\bibliofont}{\small}
\newcommand{\ie}{i.e.,\xspace}
\newcommand{\eg}{e.g.,\xspace}
\newcommand{\etc}{etc.\xspace}
\newcommand{\cf}{cf.\xspace}
\newcommand{\resp}{resp.\xspace}
\newcommand{\smaller}{\small}
\newcommand{\Subsection}{\subsection}
\newcommand{\Subsubsection}{\subsubsection}
\newcommand{\skpt}{\smallskip}
\emergencystretch=3em
\numberwithin{equation}{section}\providecommand{\og}{«\,}\providecommand{\fg}{\,»}
\input{author-preamble.tex}
\renewcommand{\proofname}{Démonstration}
\renewcommand{\refname}{Références}
\begin{document}
\begin{center}{\Large Stable item 2042}\end{center}
\paragraph{Source and attribution.}
Jean-Loup Waldspurger,\par\emph{Représentations et quasi-caractères de niveau $0$ ; endoscopie}, Journal de l’École polytechnique — Mathématiques 8 (2021), publisher version, DOI 10.5802/jep.145. \href{https://jep.centre-mersenne.org/articles/10.5802/jep.145/}{Publisher article}. Authors' text: \href{https://creativecommons.org/licenses/by/4.0/}{CC BY 4.0}.
\paragraph{Editorial problem boundary.}
Prove only the independent Proposition in Section 3.6 below, under the exact original large-residue-characteristic, facet fixed-point, compact Kottwitz-grade, cuspidal and Levi assumptions. The explicit orbital-distribution pairing, its vanishing alternatives and every Haar-volume normalisation are retained, including the conjugate on the normaliser-translated function. The entire new constant-term, root-group/segment and coset-measure core is retained. Endoscopic-transfer and later formal corollary claims are context only. Author prose remains in French.
\paragraph{Difficulty (editorial): H3.}
The definitions of parahoric cuspidal functions and Levi induction do not alone identify this orbital-distribution pairing or its exact volume factors. The author proves a constant-term calculation by facet segments and root-group cuspidal vanishing, then performs two nested coset decompositions with their Haar measures. Those geometric and measure comparisons produce the normaliser-indexed pairing, including its conjugation and the stated vanishing alternatives.
\paragraph{Editorial transcription note.}
Original author language, mathematics and ordering are preserved. Publisher front matter is replaced by this independent article wrapper. Bibliography is recovered from matching cached publisher HTML, including its TeX formulas. No new proof or mathematical repair is supplied. Surrounding original results are retained for conservative context and are not extra selected corpus claims.
\clearpage
\input{author-body.tex}
\begin{thebibliography}{99}
\bibitem{A1} J. Arthur - “The local behaviour of weighted orbital integrals” , Duke Math. J. 56 (1988) no. 2, p. 223-293
\bibitem{A2} J. Arthur - “On elliptic tempered characters” , Acta Math. 171 (1993) no. 1, p. 73-138
\bibitem{A3} J. Arthur - “On local character relations” , Selecta Math. (N.S.) 2 (1996) no. 4, p. 501-579
\bibitem{A} A.-M. Aubert - “Séries de Harish-Chandra de modules et correspondance de Howe modulaire” , J. Algebra 165 (1994) no. 3, p. 576-601
\bibitem{B} J. N. Bernstein - “Le “centre” de Bernstein” , in Representations of reductive groups over a local field (J. N. Bernstein, P. Deligne, D. Kazhdan \& M.-F. Vignéras, eds.) , Travaux en Cours , Hermann, Paris, 1984, p. 1-32, rédigé par P. Deligne
\bibitem{BT} F. Bruhat \& J. Tits - “Groupes réductifs sur un corps local : I. Données radicielles valuées” , Publ. Math. Inst. Hautes Études Sci. 41 (1972), p. 5-251
\bibitem{C} W. Casselman - “Characters and Jacquet modules” , Math. Ann. 230 (1977) no. 2, p. 101-105
\bibitem{D} S. DeBacker - “Homogeneity results for invariant distributions of a reductive $p$-adic group” , Ann. Sci. École Norm. Sup. (4) 35 (2002) no. 3, p. 391-422
\bibitem{DR} S. DeBacker \& M. Reeder - “Depth-zero supercuspidal $L$-packets and their stability” , Ann. of Math. (2) 169 (2009) no. 3, p. 795-901
\bibitem{F} A. Ferrari - “Théorème de l’indice et formule des traces” , Manuscripta Math. 124 (2007) no. 3, p. 363-390
\bibitem{HR} T. Haines \& M. Rapoport - “On parahoric subgroups” , appendice à [28]
\bibitem{HC} Harish-Chandra - Admissible invariant distributions on reductive $p$-adic groups , University Lecture Series , vol. 16 , American Mathematical Society, Providence, RI, 1999, rédigé par S. DeBacker et P. Sally
\bibitem{HV} G. Henniart \& M.-F. Vignéras - “A Satake isomorphism for representations modulo $p$ of reductive groups over local fields” , J. reine angew. Math. 701 (2015), p. 33-75
\bibitem{H} K. Hiraga - “On functoriality of Zelevinski involutions” , Compositio Math. 140 (2004) no. 6, p. 1625-1656
\bibitem{Kh} D. Kazhdan - “Cuspidal geometry of $p$-adic groups” , J. Analyse Math. 47 (1986), p. 1-36
\bibitem{KV} D. Kazhdan \& Y. Varshavsky - “Endoscopic decomposition of certain depth zero representations” , in Studies in Lie theory , Progress in Math. , vol. 243 , Birkhäuser Boston, Boston, MA, 2006, p. 223-301
\bibitem{KM} J.-L. Kim \& F. Murnaghan - “Character expansions and unrefined minimal $K$-types” , Amer. J. Math. 125 (2003) no. 6, p. 1199-1234
\bibitem{K} R. E. Kottwitz - “Isocrystals with additional structure. II” , Compositio Math. 109 (1997) no. 3, p. 255-339
\bibitem{Lanard1} T. Lanard - “Sur les $\ell $-blocs de niveau zéro des groupes $p$-adiques” , Compositio Math. 154 (2018) no. 7, p. 1473-1507
\bibitem{Lanard2} T. Lanard - “Sur les $\ell $-blocs de niveau zéro des groupes $p$-adiques II” , 2018
\bibitem{L} R. P. Langlands - “Stable conjugacy : definitions and lemmas” , Canad. J. Math. 31 (1979) no. 4, p. 700-725
\bibitem{LS} R. P. Langlands \& D. Shelstad - “On the definition of transfer factors” , Math. Ann. 278 (1987) no. 1-4, p. 219-271
\bibitem{Lu} G. Lusztig - “Classification of unipotent representations of simple $p$-adic groups” , Internat. Math. Res. Notices (1995) no. 11, p. 517-589
\bibitem{MW} C. Moeglin \& J.-L. Waldspurger - Stabilisation de la formule des traces tordue , Progress in Math. , vol. 316 \& 317 , Birkhäuser/Springer, Cham, 2016
\bibitem{MP1} A. Moy \& G. Prasad - “Unrefined minimal $K$-types for $p$-adic groups” , Invent. Math. 116 (1994) no. 1-3, p. 393-408
\bibitem{MP} A. Moy \& G. Prasad - “Jacquet functors and unrefined minimal $K$-types” , Comment. Math. Helv. 71 (1996) no. 1, p. 98-121
\bibitem{Oi} M. Oi - “Depth preserving property of the local Langlands correspondence for quasi-split classical groups in a large residual characteristic” , 2018
\bibitem{P-R08} G. Pappas \& M. Rapoport - “Twisted loop groups and their affine flag varieties” , Adv. Math. 219 (2008) no. 1, p. 118-198
\bibitem{P} G. Prasad - “Finite group actions on reductive groups and buildings and tamely-ramified descent in Bruhat-Tits theory” , Amer. J. Math. 142 (2020) no. 4, p. 1239-1267
\bibitem{PY} G. Prasad \& J.-K. Yu - “On finite group actions on reductive groups and buildings” , Invent. Math. 147 (2002) no. 3, p. 545-560
\bibitem{S} M. Solleveld - “On unipotent representations of ramified $p$-adic groups” , 2019
\bibitem{T} J. Tits - “Reductive groups over local fields” , in Automorphic forms, representations and $L$-functions (Oregon State Univ., Corvallis, Ore., 1977), Part 1 , Proc. Sympos. Pure Math. , vol. XXXIII , American Mathematical Society, Providence, RI, 1979, p. 29-69
\bibitem{W4} J.-L. Waldspurger - L’endoscopie tordue n’est pas si tordue , Mem. Amer. Math. Soc. , vol. 194 no. 908 , American Mathematical Society, Providence, RI, 2008
\bibitem{W1} J.-L. Waldspurger - “Une formule intégrale reliée à la conjecture locale de Gross-Prasad” , Compositio Math. 146 (2010) no. 5, p. 1180-1290
\bibitem{W3} J.-L. Waldspurger - “Une formule intégrale reliée à la conjecture locale de Gross-Prasad, 2 e partie : extension aux représentations tempérées” , in Sur les conjectures de Gross et Prasad. I , Astérisque , vol. 346 , Société Mathématique de France, Paris, 2012, p. 171-312
\bibitem{W2} J.-L. Waldspurger - “Caractères de représentations de niveau 0” , Ann. Fac. Sci. Toulouse Math. (6) 27 (2018) no. 5, p. 925-984
\end{thebibliography}
\end{document}
